\documentclass[11pt]{article}

\usepackage[margin=1in]{geometry}
\usepackage{amsmath, amssymb}
\usepackage{booktabs}
\usepackage{hyperref}
\usepackage{url}
\usepackage{enumitem}
\setlist{leftmargin=*}
\setlength{\parskip}{0.35em}
\setlength{\parindent}{0pt}

\title{\textbf{Stage 3: Reproduction and Implementation}\\
Implementation of Graph Partitioning Algorithms and Reproduction Results}
\author{Team Members: Name 1, Name 2, Name 3}
\date{CS240: Algorithm Design and Analysis, Spring 2026}

\begin{document}

\maketitle

\section{Implementation Overview}

The implementation is organized as a small Python project. The code is placed
under \texttt{src/}, with separate modules for graph generation, algorithms,
metrics, simulation, experiment execution, and plotting. This modular
organization makes it easier to test each component independently and to
compare several graph partitioning methods under the same evaluation pipeline.

The main files are:
\begin{itemize}
    \item \texttt{src/graphs.py}: benchmark graph construction;
    \item \texttt{src/algorithms.py}: greedy, spectral, and Kernighan--Lin
    partitioning algorithms;
    \item \texttt{src/metrics.py}: partition quality metrics;
    \item \texttt{src/simulation.py}: consensus averaging simulation;
    \item \texttt{src/experiments.py}: experiment runner and CSV export;
    \item \texttt{src/plotting.py}: generation of result figures.
\end{itemize}

The project uses Python with NumPy, NetworkX, pandas, matplotlib, and pytest.
The implementation is deterministic for the reported experiments by using a
fixed random seed.

\section{Implemented Algorithms}

\subsection{Greedy Balanced Partitioning}

The greedy baseline assigns vertices to partitions one at a time. Vertices are
processed in decreasing weighted degree order so that high-impact vertices are
placed early. For each vertex, the algorithm considers the partitions that have
not reached their capacity. It chooses the feasible partition that gives the
lowest incremental cut cost with respect to already assigned neighbors. The
capacity is set to approximately $\lceil n/k\rceil$, so the method maintains
near-perfect balance.

This baseline is intentionally simple. It provides a fast comparison point and
helps show whether more sophisticated algorithms actually improve partition
quality.

\subsection{Recursive Spectral Bisection}

The spectral method recursively bisects the largest current group until $k$
groups are obtained. For a subgraph, it constructs the weighted adjacency
matrix $A$, the degree matrix $D$, and the graph Laplacian $L=D-A$. It then
computes the eigenvectors of $L$ and uses the Fiedler vector to sort vertices.
The sorted list is split at the midpoint.

This method is useful because it gives a strong initial partition for
structured graphs. In grid and geometric graphs, the Fiedler vector often
captures spatial structure, so spectral bisection can produce low-cut balanced
partitions.

\subsection{Kernighan--Lin $k$-Way Refinement}

The implemented Kernighan--Lin method starts from a spectral partition. It then
performs several global refinement passes. During each pass, the method loops
over every pair of parts. For a pair $(a,b)$, it extracts the subgraph induced
by vertices currently assigned to either $a$ or $b$. NetworkX's robust
Kernighan--Lin bisection routine is applied to this two-part subproblem using
the current assignment as the initial partition.

The candidate update is accepted only if it does not increase the global
cut-weight objective. This acceptance rule makes the refinement monotone with
respect to the measured cut value. The refinement stops early if a full global
pass fails to improve the objective.

\section{Reproduction Scope}

The original Kernighan--Lin paper focuses on graph bisection. Our project
reproduces the core algorithmic idea rather than the exact historical
experimental setup. The reproduced components are:
\begin{enumerate}
    \item local-search refinement based on swapping vertices between two
    sides;
    \item repeated passes until no meaningful cut improvement is found;
    \item use of the method as a refinement procedure rather than a standalone
    random partitioner;
    \item empirical measurement of cut improvement and runtime cost.
\end{enumerate}

The main adaptation is the $k$-way setting. Instead of one global bisection, we
apply pairwise Kernighan--Lin refinement to pairs of parts inside a $k$-way
partition. This is a practical adaptation for distributed optimization because
the target system has $k$ workers, not only two processors.

\section{Benchmark Instances}

The implementation evaluates the algorithms on four graph families:
\begin{itemize}
    \item \textbf{Grid graphs}: structured graphs with strong local geometry;
    \item \textbf{Erdos--Renyi random graphs}: unstructured random graphs;
    \item \textbf{Random geometric graphs}: spatial graphs with local
    connectivity;
    \item \textbf{Zachary karate club graph}: a small real social network.
\end{itemize}

The benchmark suite contains eight instances:
\[
\begin{array}{lrrrr}
\toprule
\text{Graph} & \text{Size Parameter} & |V| & |E| & k \\
\midrule
\text{grid} & 6 & 36 & 60 & 4 \\
\text{grid} & 8 & 64 & 112 & 4 \\
\text{grid} & 10 & 100 & 180 & 4 \\
\text{erdos\_renyi} & 40 & 40 & 93 & 4 \\
\text{erdos\_renyi} & 80 & 80 & 248 & 4 \\
\text{geometric} & 50 & 50 & 330 & 4 \\
\text{geometric} & 90 & 90 & 628 & 4 \\
\text{karate} & 34 & 34 & 78 & 2 \\
\bottomrule
\end{array}
\]

\section{Metrics}

The primary metric is cut weight:
\[
    \mathrm{cut}(p)=\sum_{(u,v)\in E,\;p(u)\neq p(v)}w_{uv}.
\]
Lower cut weight means fewer or cheaper cross-worker dependencies.

The second metric is normalized cut, which divides boundary weight by the
volume of each part and sums the ratios. This helps evaluate whether the cut is
small relative to the connectivity of each partition.

The third metric is balance ratio:
\[
    \rho(p)=
    \frac{\max_i |\{v:p(v)=i\}|}{|V|/k}.
\]
Values close to $1$ indicate balanced workload.

The fourth metric is runtime in seconds. This measures the computational cost
of each partitioning algorithm.

Finally, the consensus simulation reports final consensus error and the number
of cross-partition messages. The simulation keeps the original communication
graph fixed, so partitions do not change the convergence dynamics; instead,
they change how many messages cross worker boundaries.

\section{Reproduced Results}

The experiments reproduce the expected qualitative behavior of
Kernighan--Lin-style refinement. Across the benchmark suite, the average cut
weights are:
\[
    \text{greedy}=98.25,\qquad
    \text{spectral}=72.75,\qquad
    \text{Kernighan--Lin}=66.88.
\]
The average cross-partition message counts in the consensus simulation are:
\[
    \text{greedy}=15460,\qquad
    \text{spectral}=11360,\qquad
    \text{Kernighan--Lin}=10440.
\]

The Kernighan--Lin refinement achieves the lowest average cut and the lowest
average cross-worker traffic. It gives the best cut on six of the eight graph
instances. This supports the intended reproduction target: local refinement
often improves partition quality compared with a simpler baseline and can also
improve a spectral initialization.

\section{Deviation and Failure Analysis}

There are two important deviations from the original paper. First, the
original Kernighan--Lin method is a two-way graph partitioning algorithm, while
our application requires $k$ parts. We therefore use pairwise two-way
refinement as a component inside a $k$-way framework.

Second, the project uses NetworkX's implementation of pairwise
Kernighan--Lin bisection rather than reimplementing every low-level detail
from scratch. This choice improves robustness and lets the project focus on
algorithmic comparison, adaptation, and evaluation. The core reproduced
algorithmic behavior remains the local-search cut refinement.

The consensus experiment also has a limitation: it measures cross-worker
traffic but does not alter the underlying communication topology. Therefore,
different partitions have nearly identical consensus error on the same graph.
This is not a failure of the partitioning algorithms; it reflects the design of
the simulation. The simulation isolates communication placement cost rather
than changing the numerical averaging process.

\section{How to Run}

The experiments can be reproduced with:
\[
\texttt{python src/experiments.py --output-dir outputs --seed 7}
\]
This command writes the raw CSV file to
\texttt{outputs/partition\_results.csv} and saves figures under
\texttt{outputs/figures/}. Unit tests can be run with:
\[
\texttt{python -m pytest}
\]

\end{document}
